\documentclass[11pt,a4paper]{article}
\usepackage[margin=1in]{geometry}
\usepackage[T1]{fontenc}
\usepackage[utf8]{inputenc}
\usepackage{lmodern}
\usepackage{hyperref}
\usepackage{longtable}
\usepackage{booktabs}
\usepackage{array}
\usepackage{enumitem}
\usepackage{xcolor}
\usepackage{listings}
\usepackage{titlesec}
\n\hypersetup{colorlinks=true, linkcolor=blue, urlcolor=blue}
\n\lstset{
  language=Python,
  basicstyle=\ttfamily\scriptsize,
  breaklines=true,
  frame=single,
  tabsize=4,
  showstringspaces=false,
  columns=fullflexible
}
\n\title{Technical Documentation: \texttt{basic\_pipelines/human\_group\_tracking.py}}
\author{Repository: \texttt{hailo-rpi5-examples}}
\date{\today}
\n\begin{document}
\maketitle
\tableofcontents
\newpage
\n\section{Document Scope}
This document provides a code-centric technical description of \texttt{basic\_pipelines/human\_group\_tracking.py}, including:
\begin{itemize}[leftmargin=*]
  \item Hailo inference pipeline integration.
  \item Group clustering and lock-on targeting logic.
  \item PCA9685 servo maneuver control logic.
  \item Model artifact usage (HEF/labels) as implemented in this codebase.
  \item Full source code inclusion.
\end{itemize}

All statements in this document are derived from files present in the repository.

\section{Codebase Evidence Set}
The following files were used to build this documentation:
\begin{itemize}[leftmargin=*]
  \item \texttt{basic\_pipelines/human\_group\_tracking.py}
  \item \texttt{basic\_pipelines/human\_group\_tracking\_backup.py}
  \item \texttt{basic\_pipelines/main.py}
  \item \texttt{local\_resources/best\_person.json}
  \item \texttt{.env}
  \item \texttt{README.md}
\end{itemize}

\section{System-Level Overview}
\subsection{Runtime Flow}
\begin{enumerate}[leftmargin=*]
  \item Input video frames are fed to a Hailo GStreamer detection pipeline.
  \item Detection metadata is extracted in \texttt{app\_callback}.
  \item Person detections are clustered by spatial proximity.
  \item A target group is selected (largest group, then lock-follow behavior).
  \item Pan/Tilt offsets are computed relative to frame center.
  \item Servo controller applies proportional camera adjustments over PCA9685.
\end{enumerate}

\subsection{Primary Module Coupling}
\begin{longtable}{>{\raggedright\arraybackslash}p{0.28\textwidth} p{0.64\textwidth}}
\toprule
\textbf{Component} & \textbf{Function} \\
\midrule
\texttt{GStreamerDetectionApp} & Executes Hailo detection pipeline and callback loop. \\
\texttt{app\_callback} & Parses detections, clusters humans, determines target group, overlays visuals, updates servos. \\
\texttt{ServoController} & Initializes PCA9685 and applies pan/tilt/focus angle commands. \\
\texttt{KeyboardListener} & Captures terminal key commands for focus adjustments and quit. \\
\bottomrule
\end{longtable}

\section{Model Artifacts and Inference Inputs}
\subsection{Configured Defaults}
From \texttt{human\_group\_tracking.py}:
\begin{itemize}[leftmargin=*]
  \item Default HEF path: \texttt{/usr/local/hailo/resources/models/hailo8l/best.hef}
  \item Default labels JSON: \texttt{local\_resources/best\_person.json}
\end{itemize}

From \texttt{local\_resources/best\_person.json}:
\begin{itemize}[leftmargin=*]
  \item \texttt{iou\_threshold = 0.45}
  \item \texttt{detection\_threshold = 0.3}
  \item Label set contains one class: \texttt{person}
\end{itemize}

\subsection{Class Filtering Behavior}
The callback currently treats all detections as persons in \texttt{human\_group\_tracking.py}. The backup variant includes explicit permissive single-class logic (label \texttt{person}, user target label, or class ID 0), indicating support for custom single-class models.

\section{Tracking and Group Selection Logic}
\subsection{Clustering}
The active implementation uses \texttt{simple\_cluster()} with union-find style transitive merging and Euclidean distance threshold \texttt{eps}. Cluster labels are generated without external ML clustering dependencies.

\subsection{Target Selection}
\begin{enumerate}[leftmargin=*]
  \item Compute cluster-level bounding boxes and centers for clusters with at least 2 members.
  \item Prefer largest valid cluster for initial target.
  \item When lock is active, follow nearest cluster to prior lock center within lock radius.
  \item On lock loss, auto-retarget to largest available group.
\end{enumerate}

\subsection{Pan/Tilt Offset Computation}
\begin{itemize}[leftmargin=*]
  \item Frame center: \((W/2, H/2)\)
  \item Target center: group center \((x_t, y_t)\)
  \item Pixel offsets: \(\Delta x = x_t - W/2\), \(\Delta y = y_t - H/2\)
  \item Percentage offsets used for control: \(100\Delta x/W\), \(100\Delta y/H\)
\end{itemize}

\section{Servo Maneuver Logic}
\subsection{Hardware and Channels}
In the active script:
\begin{itemize}[leftmargin=*]
  \item Pan channel constant: 1
  \item Tilt channel constant: 0
  \item Focus channel constant: 2
\end{itemize}
At runtime, CLI defaults are pan=0, tilt=1, focus=2 and can be swapped/inverted by arguments.

\subsection{Servo Control Parameters}
\begin{longtable}{>{\ttfamily\raggedright\arraybackslash}p{0.28\textwidth} p{0.18\textwidth} p{0.46\textwidth}}
\toprule
\textbf{Parameter} & \textbf{Value} & \textbf{Role} \\
\midrule
SERVO\_FREQ & 50 & PWM frequency for servo drive. \\
SERVO\_MIN\_PULSE & 500 & Minimum pulse width (\textmu s). \\
SERVO\_MAX\_PULSE & 2500 & Maximum pulse width (\textmu s). \\
DEAD\_ZONE & 12 & Offset threshold (\%) below which movement is suppressed. \\
UPDATE\_INTERVAL & 0.25 & Rate limit between updates (seconds). \\
base\_step & 0.8 & Minimum movement step in degrees. \\
max\_step & 3.0 & Maximum movement step in degrees. \\
\bottomrule
\end{longtable}

\subsection{Control Law}
Pan and tilt commands are proportional to target offset magnitude with clamped step ranges. Additional options support inversion and channel swapping for installation-specific mounting orientation.

\subsection{Focus Logic}
Manual focus control increments/decrements focus angle by 5 degrees per key event, clamped to [0,360].

\section{CLI and Operational Interface}
\subsection{Arguments in \texttt{human\_group\_tracking.py}}
\begin{itemize}[leftmargin=*]
  \item \texttt{--clustering-eps}
  \item \texttt{--no-servo}
  \item \texttt{--pan-channel}
  \item \texttt{--tilt-channel}
  \item \texttt{--focus-channel}
  \item \texttt{--servo-speed}
  \item \texttt{--invert-pan}
  \item \texttt{--invert-tilt}
  \item \texttt{--swap-servos}
\end{itemize}

\subsection{Keyboard Inputs}
\begin{itemize}[leftmargin=*]
  \item In OpenCV window: \texttt{s} retarget, up/down arrows or \texttt{w/x} for focus.
  \item In terminal listener thread: \texttt{w/x} focus, \texttt{s} retarget signal, \texttt{q} immediate exit.
\end{itemize}

\section{Training and HEF Conversion Process in This Repository}
\subsection{What Is Explicitly Present}
\begin{itemize}[leftmargin=*]
  \item \texttt{basic\_pipelines/main.py} exposes modes \texttt{train}, \texttt{test}, \texttt{webcam}, and \texttt{hailo}.
  \item The \texttt{hailo} mode executes \texttt{human\_group\_tracking.py}.
  \item The tracking script consumes a precompiled HEF artifact (default: \texttt{best.hef}).
  \item The backup script explicitly describes the model as a custom trained model for tiny persons.
\end{itemize}

\subsection{What Is Not Present in Current Tree}
\begin{itemize}[leftmargin=*]
  \item \texttt{basic\_pipelines/src/train.py} and \texttt{basic\_pipelines/src/test.py} are referenced by launcher code but not present in this repository snapshot.
  \item No in-repo script documents dataset preparation, tiny-people training hyperparameters, or Hailo Dataflow Compiler command sequence producing \texttt{best.hef}.
\end{itemize}

\subsection{Implication for Reproducibility}
The runtime inference and tracking stack are fully documented by this repository, but the exact model training and HEF compilation pipeline are external to this code tree.

\section{Environment and Version Signals}
From \texttt{.env}:
\begin{itemize}[leftmargin=*]
  \item \texttt{host\_arch=rpi}
  \item \texttt{hailo\_arch=hailo8l}
  \item \texttt{model\_zoo\_version=v2.14.0}
  \item \texttt{hailort\_version=4.23.0}
  \item \texttt{tappas\_version=5.1.0}
\end{itemize}

\section{Execution Sequence}
\begin{lstlisting}[language=bash]
source setup_env.sh
python basic_pipelines/human_group_tracking.py --input rpi
# Optional:
#   --no-servo
#   --clustering-eps 120
#   --pan-channel 0 --tilt-channel 1 --focus-channel 2
#   --invert-pan --invert-tilt --swap-servos
\end{lstlisting}

\section{Failure Modes and Mitigations in Code}
\begin{itemize}[leftmargin=*]
  \item PCA9685 init retries (3 attempts) with delay.
  \item I2C write retries in \texttt{\_safe\_set\_angle}.
  \item Dead-zone suppression to reduce servo chatter.
  \item Update rate limiting to avoid oscillatory overdriving.
  \item Lock-loss retargeting fallback.
\end{itemize}

\section{Technical Gaps to Fill for End-to-End Model Lifecycle}
For complete training-to-HEF traceability, add or link the missing artifacts:
\begin{itemize}[leftmargin=*]
  \item Training script and exact CLI used for tiny-people dataset.
  \item Dataset spec (splits, annotation format, preprocessing).
  \item Quantization/optimization and DFC compilation commands.
  \item Validation metrics tied to \texttt{best.hef} build.
\end{itemize}

\newpage
\section{Full Source Code: \texttt{basic\_pipelines/human\_group\_tracking.py}}
\begin{lstlisting}
"""
Human Group Tracking - With Servo Control
Detects humans, clusters by proximity, automatically adjusts pan/tilt servos
to center the largest group in the frame.

Hardware:
    - Pan servo on GPIO 23
    - Tilt servo on GPIO 24

Usage:
    python human_group_tracking.py --input rpi
    python human_group_tracking.py --input rpi --no-servo  # Disable servo control
"""

from pathlib import Path
import gi
gi.require_version('Gst', '1.0')
from gi.repository import Gst
import os
import sys
import hailo
import argparse
import cv2
import numpy as np
import time
import threading
import select

from hailo_apps.hailo_app_python.core.common.buffer_utils import get_caps_from_pad, get_numpy_from_buffer
from hailo_apps.hailo_app_python.core.gstreamer.gstreamer_app import app_callback_class
from hailo_apps.hailo_app_python.apps.detection.detection_pipeline import GStreamerDetectionApp

# Default paths
DEFAULT_CUSTOM_HEF = "/usr/local/hailo/resources/models/hailo8l/best.hef"
DEFAULT_LABELS_JSON = str(Path(__file__).resolve().parent.parent / "local_resources" / "best_person.json")

# Servo PCA9685 channels (instead of GPIO pins)
PAN_CHANNEL = 1  # PCA9685 channel for pan servo
TILT_CHANNEL = 0  # PCA9685 channel for tilt servo
FOCUS_CHANNEL = 2  # PCA9685 channel for focus servo (manual control)

# Servo parameters
SERVO_FREQ = 50  # 50Hz for standard servos
SERVO_MIN_PULSE = 500   # microseconds (0 degrees)
SERVO_MAX_PULSE = 2500  # microseconds (180 degrees)
SERVO_CENTER = 1500     # microseconds (90 degrees / center)
SERVO_INIT = 1500       # Initialize at 90 degrees (center)

# Control parameters
DEAD_ZONE = 12          # Percentage - don't move if offset is less than this (was 8)
SPEED_FACTOR = 0.8    # How fast servos respond (0.1 = slow, 1.0 = fast)
UPDATE_INTERVAL = 0.25  # Seconds between servo updates (was 0.2)
SMOOTHING = 0.7         # Exponential smoothing (0.9 = very smooth, 0.5 = responsive)
MIN_PULSE_CHANGE = 15   # Minimum pulse change (us) to actually move servo

# -----------------------------------------------------------------------------------------------
# Servo Controller using PCA9685 (I2C PWM driver - very stable)
# -----------------------------------------------------------------------------------------------
class ServoController:
    def __init__(self, pan_channel=PAN_CHANNEL, tilt_channel=TILT_CHANNEL, focus_channel=FOCUS_CHANNEL,
                 speed=SPEED_FACTOR, invert_pan=False, invert_tilt=False, enabled=True):
        self.pan_channel = pan_channel
        self.tilt_channel = tilt_channel
        self.focus_channel = focus_channel
        self.speed = speed
        self.invert_pan = invert_pan
        self.invert_tilt = invert_tilt
        self.enabled = enabled
        self.pca = None
        self.pan_servo = None
        self.tilt_servo = None
        self.focus_servo = None
        
        # Current positions as angle (0-180 degrees, 90 = center)
        self.pan_angle = 90.0  # Start at center
        self.tilt_angle = 90.0
        self.focus_angle = 180.0  # Focus starts at center (0-360 range)
        
        # Target positions (smoothed)
        self.pan_target = 90.0
        self.tilt_target = 90.0
        
        # Last sent positions (to avoid redundant updates)
        self.last_pan_sent = 90.0
        self.last_tilt_sent = 90.0
        
        # Last update time
        self.last_update = 0
        
        # Track if we have active detection
        self.has_target = False
        
        if self.enabled:
            self._init_pca9685()
    
    def _init_pca9685(self):
        """Initialize PCA9685 servo driver"""
        try:
            import board
            import busio
            from adafruit_pca9685 import PCA9685
            from adafruit_motor import servo
            
            # Initialize I2C with explicit frequency (lower = more stable)
            # Default is 100kHz, some PCA9685 boards have issues at high speeds
            i2c = busio.I2C(board.SCL, board.SDA, frequency=100000)
            
            # Retry logic for PCA9685 initialization
            for attempt in range(3):
                try:
                    self.pca = PCA9685(i2c)
                    self.pca.frequency = SERVO_FREQ
                    break
                except Exception as e:
                    print(f"PCA9685 init attempt {attempt+1} failed: {e}")
                    time.sleep(0.5)
                    if attempt == 2:
                        raise
            
            # Create servo objects on specified channels
            # actuation_range is the total range in degrees
            # min_pulse and max_pulse in microseconds
            self.pan_servo = servo.Servo(
                self.pca.channels[self.pan_channel],
                min_pulse=SERVO_MIN_PULSE,
                max_pulse=SERVO_MAX_PULSE,
                actuation_range=180
            )
            self.tilt_servo = servo.Servo(
                self.pca.channels[self.tilt_channel],
                min_pulse=SERVO_MIN_PULSE,
                max_pulse=SERVO_MAX_PULSE,
                actuation_range=180
            )
            self.focus_servo = servo.Servo(
                self.pca.channels[self.focus_channel],
                min_pulse=SERVO_MIN_PULSE,
                max_pulse=SERVO_MAX_PULSE,
                actuation_range=360  # 360 degree range for focus
            )
            
            # Move to center with retry
            self._safe_set_angle(self.pan_servo, 90)   # Pan center = 90
            self._safe_set_angle(self.tilt_servo, 90)  # Tilt center = 90
            self._safe_set_angle(self.focus_servo, 180) # Focus center = 180 (middle of 0-360)
            time.sleep(0.5)  # Give servos time to reach position
            
            print(f"PCA9685 Servos initialized: Pan=CH{self.pan_channel}, Tilt=CH{self.tilt_channel}, Focus=CH{self.focus_channel}")
            print(f"Servos at 90 degrees (center)")
            
        except Exception as e:
            print(f"Warning: Could not initialize PCA9685: {e}")
            import traceback
            traceback.print_exc()
            self.enabled = False
            self.pca = None
            self.pan_servo = None
            self.tilt_servo = None
    
    def _safe_set_angle(self, servo_obj, angle, retries=2):
        """Set servo angle with retry on I2C errors"""
        if servo_obj is None:
            return False
        
        angle = max(0, min(180, angle))
        
        for attempt in range(retries + 1):
            try:
                servo_obj.angle = angle
                return True
            except OSError as e:
                if attempt < retries:
                    time.sleep(0.01)  # Small delay before retry
                # Silently fail after retries - don't spam console
        return False
    
    def _set_servo_angle(self, servo_obj, angle):
        """Set servo position using angle (0-180 degrees)"""
        if not self.enabled or servo_obj is None:
            return
        self._safe_set_angle(servo_obj, angle)
    
    def _release_servos(self):
        """Release servos - keep at current position, just stop updating"""
        # Don't set angle to None as this can cause I2C issues
        # Just stop calling update - servos will hold their last position
        pass
    
    def update(self, pan_offset_pct, tilt_offset_pct, has_detection=True):
        """
        Update servo positions to CENTER the target in the frame.
        pan_offset_pct: positive = target is to the RIGHT of frame center
        tilt_offset_pct: positive = target is BELOW frame center
        has_detection: if False, servos hold their current position
        
        Logic:
        - Target on RIGHT of frame -> camera must rotate RIGHT -> target moves LEFT toward center
        - Target on LEFT of frame -> camera must rotate LEFT -> target moves RIGHT toward center
        - Target BELOW center -> camera must tilt DOWN -> target moves UP toward center
        - Target ABOVE center -> camera must tilt UP -> target moves DOWN toward center
        """
        if not self.enabled:
            return
        
        self.has_target = has_detection
        
        # If no detection, hold current position
        if not has_detection:
            return
        
        # Rate limiting
        now = time.time()
        if now - self.last_update < UPDATE_INTERVAL:
            return
        self.last_update = now
        
        # Store raw offsets for logging
        raw_pan_offset = pan_offset_pct
        raw_tilt_offset = tilt_offset_pct
        
        # Dead zone - target is in center band, stop moving
        if abs(pan_offset_pct) < DEAD_ZONE:
            pan_offset_pct = 0
        if abs(tilt_offset_pct) < DEAD_ZONE:
            tilt_offset_pct = 0
        
        # If target is centered, hold position
        if pan_offset_pct == 0 and tilt_offset_pct == 0:
            print(f"[SERVO] Target in dead zone (pan={raw_pan_offset:+.1f}%, tilt={raw_tilt_offset:+.1f}%) - HOLDING")
            return
        
        # PROPORTIONAL step size - move faster when far, slower when close
        # This prevents overshoot by slowing down as we approach the target
        base_step = 0.8  # Minimum step in degrees
        max_step = 3.0   # Maximum step in degrees
        
        # Scale step by how far off-center we are (0-50% range mapped to base-max)
        pan_magnitude = min(abs(pan_offset_pct), 50) / 50.0  # 0.0 to 1.0
        tilt_magnitude = min(abs(tilt_offset_pct), 50) / 50.0
        
        pan_step = base_step + (max_step - base_step) * pan_magnitude
        tilt_step = base_step + (max_step - base_step) * tilt_magnitude
        
        # Determine direction
        pan_adjust = 0.0
        tilt_adjust = 0.0
        
        # PAN: Target on RIGHT (+) -> need to turn camera RIGHT to bring target to center
        #      Assuming: increasing angle = camera turns RIGHT
        if pan_offset_pct > 0:
            pan_adjust = pan_step
        elif pan_offset_pct < 0:
            pan_adjust = -pan_step
        
        # TILT: Target BELOW (+) -> need to tilt camera DOWN to bring target to center
        #       Assuming: increasing angle = camera tilts DOWN
        if tilt_offset_pct > 0:
            tilt_adjust = tilt_step
        elif tilt_offset_pct < 0:
            tilt_adjust = -tilt_step
        
        # Apply inversion if servo is mounted opposite
        if self.invert_pan:
            pan_adjust = -pan_adjust
        if self.invert_tilt:
            tilt_adjust = -tilt_adjust
        
        # Store old angles for comparison
        old_pan = self.pan_angle
        old_tilt = self.tilt_angle
        
        # Update angle
        self.pan_angle = self.pan_angle + pan_adjust
        self.tilt_angle = self.tilt_angle + tilt_adjust
        
        # Clamp to valid range (0-180 degrees)
        self.pan_angle = max(0, min(180, self.pan_angle))
        self.tilt_angle = max(0, min(180, self.tilt_angle))
        
        # Simplified logging
        print(f"[SERVO] Target: pan={raw_pan_offset:+.1f}% tilt={raw_tilt_offset:+.1f}% | Step: pan={pan_adjust:+.1f}° tilt={tilt_adjust:+.1f}° | Angle: {old_pan:.0f}°->{self.pan_angle:.0f}° / {old_tilt:.0f}°->{self.tilt_angle:.0f}°")
        
        # Send to servos
        self._set_servo_angle(self.pan_servo, self.pan_angle)
        self._set_servo_angle(self.tilt_servo, self.tilt_angle)
        self.last_pan_sent = self.pan_angle
        self.last_tilt_sent = self.tilt_angle
    
    def adjust_focus(self, direction):
        """
        Manually adjust focus servo.
        direction: +1 for increase (up arrow), -1 for decrease (down arrow)
        """
        if not self.enabled or self.focus_servo is None:
            return
        
        step = 5.0  # Degrees per key press (larger for 360 range)
        old_focus = self.focus_angle
        self.focus_angle = self.focus_angle + (direction * step)
        self.focus_angle = max(0, min(360, self.focus_angle))  # 0-360 range
        
        self._set_servo_angle(self.focus_servo, self.focus_angle)
        print(f"[FOCUS] {old_focus:.0f}° -> {self.focus_angle:.0f}° ({'UP' if direction > 0 else 'DOWN'})")
    
    def center(self):
        """Move servos to center position"""
        self.pan_angle = 90.0
        self.tilt_angle = 90.0
        self.pan_target = 90.0
        self.tilt_target = 90.0
        if self.pan_servo:
            self.pan_servo.angle = 90
        if self.tilt_servo:
            self.tilt_servo.angle = 90
    
    def cleanup(self):
        """Cleanup PCA9685 resources"""
        try:
            # Release servos
            if self.pan_servo:
                self.pan_servo.angle = None
            if self.tilt_servo:
                self.tilt_servo.angle = None
            if self.focus_servo:
                self.focus_servo.angle = None
            # Deinit PCA9685
            if self.pca:
                self.pca.deinit()
            print("PCA9685 servos cleaned up")
        except:
            pass

# Global servo controller (initialized in main)
servo_controller = None

# -----------------------------------------------------------------------------------------------
# Simple distance-based clustering (no sklearn needed)
# -----------------------------------------------------------------------------------------------
def simple_cluster(centers, eps):
    """
    Simple clustering: group points within eps distance of each other.
    Uses union-find approach for transitive clustering.
    """
    n = len(centers)
    if n == 0:
        return []
    
    # Initialize each point as its own cluster
    parent = list(range(n))
    
    def find(x):
        if parent[x] != x:
            parent[x] = find(parent[x])
        return parent[x]
    
    def union(x, y):
        px, py = find(x), find(y)
        if px != py:
            parent[px] = py
    
    # Union points that are close enough
    for i in range(n):
        for j in range(i + 1, n):
            dx = centers[i][0] - centers[j][0]
            dy = centers[i][1] - centers[j][1]
            dist = (dx * dx + dy * dy) ** 0.5
            if dist < eps:
                union(i, j)
    
    # Convert to cluster labels
    root_to_label = {}
    labels = []
    next_label = 0
    
    for i in range(n):
        root = find(i)
        if root not in root_to_label:
            root_to_label[root] = next_label
            next_label += 1
        labels.append(root_to_label[root])
    
    return labels

# -----------------------------------------------------------------------------------------------
# User data class
# -----------------------------------------------------------------------------------------------
class user_app_callback_class(app_callback_class):
    def __init__(self, eps=100):
        super().__init__()
        self.eps = eps
        self.frame_count = 0
        
        # Lock-on tracking state
        self.locked_group_center = None  # (cx, cy) of locked group
        self.lock_active = False
        self.lock_search_radius = 150  # Pixels - how far to search for the locked group
        self.frame_height = 720  # Will be updated

# Focus slider dimensions (will be set based on frame size)
FOCUS_SLIDER_X = 20
FOCUS_SLIDER_Y = 50  # From bottom
FOCUS_SLIDER_WIDTH = 200
FOCUS_SLIDER_HEIGHT = 25

def draw_focus_slider(frame, focus_angle, height):
    """Draw a focus control slider bar on the frame"""
    if servo_controller is None or not servo_controller.enabled:
        return
    
    # Slider position (bottom left)
    x = FOCUS_SLIDER_X
    y = height - FOCUS_SLIDER_Y
    w = FOCUS_SLIDER_WIDTH
    h = FOCUS_SLIDER_HEIGHT
    
    # Background bar (dark gray)
    cv2.rectangle(frame, (x, y), (x + w, y + h), (50, 50, 50), -1)
    cv2.rectangle(frame, (x, y), (x + w, y + h), (150, 150, 150), 2)
    
    # Fill based on focus angle (0-360)
    fill_width = int((focus_angle / 360.0) * (w - 4))
    cv2.rectangle(frame, (x + 2, y + 2), (x + 2 + fill_width, y + h - 2), (0, 200, 255), -1)
    
    # Slider knob
    knob_x = x + 2 + fill_width
    cv2.rectangle(frame, (knob_x - 3, y), (knob_x + 3, y + h), (255, 255, 255), -1)
    
    # Label
    cv2.putText(frame, f"FOCUS: {focus_angle:.0f}", (x, y - 8), 
                cv2.FONT_HERSHEY_SIMPLEX, 0.5, (0, 200, 255), 1)
    cv2.putText(frame, "0", (x, y + h + 15), cv2.FONT_HERSHEY_SIMPLEX, 0.4, (150, 150, 150), 1)
    cv2.putText(frame, "360", (x + w - 25, y + h + 15), cv2.FONT_HERSHEY_SIMPLEX, 0.4, (150, 150, 150), 1)
    
    # Instructions
    cv2.putText(frame, "W/X keys = focus", (x + w + 10, y + 15), 
                cv2.FONT_HERSHEY_SIMPLEX, 0.4, (150, 150, 150), 1)

# -----------------------------------------------------------------------------------------------
# Callback - draws clusters, centers, and pan/tilt on frame
# -----------------------------------------------------------------------------------------------
def app_callback(pad, info, user_data):
    buffer = info.get_buffer()
    if buffer is None:
        return Gst.PadProbeReturn.OK
    
    user_data.frame_count += 1
    
    # Get frame dimensions
    fmt, width, height = get_caps_from_pad(pad)
    if width is None:
        width, height = 1280, 720
    
    # Store height for mouse callback
    user_data.frame_height = height
    
    # Get frame for drawing (only if use_frame is enabled)
    frame = None
    if user_data.use_frame and fmt is not None:
        frame = get_numpy_from_buffer(buffer, fmt, width, height)
    
    # Get detections
    roi = hailo.get_roi_from_buffer(buffer)
    detections = roi.get_objects_typed(hailo.HAILO_DETECTION)
    
    # All detections are persons (single class model)
    persons = list(detections)
    n_persons = len(persons)
    
    # Frame center
    frame_cx, frame_cy = width // 2, height // 2
    
    # Draw frame center marker if we have a frame
    if frame is not None:
        # Draw "+" at frame center (white with black outline)
        cv2.line(frame, (frame_cx - 20, frame_cy), (frame_cx + 20, frame_cy), (0, 0, 0), 4)
        cv2.line(frame, (frame_cx, frame_cy - 20), (frame_cx, frame_cy + 20), (0, 0, 0), 4)
        cv2.line(frame, (frame_cx - 20, frame_cy), (frame_cx + 20, frame_cy), (255, 255, 255), 2)
        cv2.line(frame, (frame_cx, frame_cy - 20), (frame_cx, frame_cy + 20), (255, 255, 255), 2)
    
    if n_persons == 0:
        if frame is not None:
            cv2.putText(frame, "No detections", (10, 30), cv2.FONT_HERSHEY_SIMPLEX, 0.8, (0, 255, 255), 2)
            frame = cv2.cvtColor(frame, cv2.COLOR_RGB2BGR)
            user_data.set_frame(frame)
        return Gst.PadProbeReturn.OK
    
    # Get centers in pixels for clustering
    centers = []
    for det in persons:
        bbox = det.get_bbox()
        cx = int((bbox.xmin() + bbox.xmax()) / 2 * width)
        cy = int((bbox.ymin() + bbox.ymax()) / 2 * height)
        centers.append((cx, cy))
    
    # Cluster persons
    labels = simple_cluster(centers, user_data.eps)
    
    # Group by cluster
    clusters = {}
    for i, lbl in enumerate(labels):
        if lbl not in clusters:
            clusters[lbl] = []
        clusters[lbl].append(i)
    
    # Find largest cluster (with 2+ members)
    largest_id = -1
    largest_size = 0
    for cid, members in clusters.items():
        if len(members) >= 2 and len(members) > largest_size:
            largest_size = len(members)
            largest_id = cid
    
    # Calculate all group centers
    group_centers = {}  # cid -> (cx, cy)
    for cid, members in clusters.items():
        if len(members) < 2:
            continue
        min_x = min_y = float('inf')
        max_x = max_y = 0
        for idx in members:
            bbox = persons[idx].get_bbox()
            min_x = min(min_x, int(bbox.xmin() * width))
            min_y = min(min_y, int(bbox.ymin() * height))
            max_x = max(max_x, int(bbox.xmax() * width))
            max_y = max(max_y, int(bbox.ymax() * height))
        group_centers[cid] = ((min_x + max_x) // 2, (min_y + max_y) // 2, len(members))
    
    # Lock-on logic: Find group to track
    target_group_id = None
    
    if user_data.lock_active and user_data.locked_group_center is not None:
        # Find group closest to last locked position
        best_dist = float('inf')
        best_cid = None
        for cid, (cx, cy, size) in group_centers.items():
            dx = cx - user_data.locked_group_center[0]
            dy = cy - user_data.locked_group_center[1]
            dist = (dx*dx + dy*dy) ** 0.5
            if dist < best_dist and dist < user_data.lock_search_radius:
                best_dist = dist
                best_cid = cid
        
        if best_cid is not None:
            target_group_id = best_cid
            # Update locked position to follow the group
            user_data.locked_group_center = (group_centers[best_cid][0], group_centers[best_cid][1])
        else:
            # Lost the locked group - auto-retarget to largest group
            print("Lock lost - auto-retargeting to largest group")
            if largest_id >= 0 and largest_id in group_centers:
                target_group_id = largest_id
                user_data.locked_group_center = (group_centers[largest_id][0], group_centers[largest_id][1])
                print(f"Retargeted to group size: {group_centers[largest_id][2]}")
            else:
                # No groups available
                target_group_id = None
                user_data.lock_active = False
                user_data.locked_group_center = None
    else:
        # Not locked or no lock - target largest group
        if largest_id >= 0:
            target_group_id = largest_id
            # Auto-lock on first detection
            if not user_data.lock_active and largest_id in group_centers:
                user_data.lock_active = True
                user_data.locked_group_center = (group_centers[largest_id][0], group_centers[largest_id][1])
    
    # Colors for clusters
    colors = [(255, 0, 0), (0, 255, 0), (0, 0, 255), (255, 255, 0), (255, 0, 255), (0, 255, 255)]
    
    target_group_center = None
    
    if frame is not None:
        # Draw each cluster
        for cid, members in clusters.items():
            if len(members) < 2:
                continue
            
            color = colors[cid % len(colors)]
            is_target = (cid == target_group_id)
            is_largest = (cid == largest_id)
            
            # Calculate cluster bounding box
            min_x = min_y = float('inf')
            max_x = max_y = 0
            
            for idx in members:
                bbox = persons[idx].get_bbox()
                min_x = min(min_x, int(bbox.xmin() * width))
                min_y = min(min_y, int(bbox.ymin() * height))
                max_x = max(max_x, int(bbox.xmax() * width))
                max_y = max(max_y, int(bbox.ymax() * height))
            
            # Add padding
            pad = 15
            min_x = max(0, min_x - pad)
            min_y = max(0, min_y - pad)
            max_x = min(width, max_x + pad)
            max_y = min(height, max_y + pad)
            
            # Draw cluster box - target gets special color (cyan) and thickness
            if is_target:
                box_color = (0, 255, 255)  # Cyan for locked target
                thickness = 4
            else:
                box_color = color
                thickness = 2
            cv2.rectangle(frame, (min_x, min_y), (max_x, max_y), box_color, thickness)
            
            # Cluster center
            group_cx = (min_x + max_x) // 2
            group_cy = (min_y + max_y) // 2
            
            # Draw "+" at group center
            size = 15 if is_target else 10
            marker_color = (0, 255, 255) if is_target else color
            cv2.line(frame, (group_cx - size, group_cy), (group_cx + size, group_cy), marker_color, 3)
            cv2.line(frame, (group_cx, group_cy - size), (group_cx, group_cy + size), marker_color, 3)
            
            # Label - show LOCKED for target, largest for biggest
            if is_target:
                label = f"LOCKED ({len(members)})"
            elif is_largest:
                label = f"GROUP {len(members)}"
            else:
                label = f"grp {len(members)}"
            cv2.putText(frame, label, (min_x, min_y - 5), cv2.FONT_HERSHEY_SIMPLEX, 0.7, box_color, 2)
            
            if is_target:
                target_group_center = (group_cx, group_cy)
        
        # Draw line from target group center to frame center
        if target_group_center is not None:
            cv2.line(frame, target_group_center, (frame_cx, frame_cy), (0, 255, 255), 2)
            
            # Calculate pan/tilt percentage
            pan_px = target_group_center[0] - frame_cx
            tilt_px = target_group_center[1] - frame_cy
            pan_pct = (pan_px / width) * 100
            tilt_pct = (tilt_px / height) * 100
            
            # Update servos to track the group
            if servo_controller is not None:
                servo_controller.update(pan_pct, tilt_pct, has_detection=True)
            
            pan_dir = "R" if pan_px > 0 else "L"
            tilt_dir = "D" if tilt_px > 0 else "U"
            
            # Draw pan/tilt info at top
            lock_status = "LOCKED" if user_data.lock_active else "AUTO"
            servo_status = "TRACKING" if servo_controller and servo_controller.enabled else "NO SERVO"
            info_text = f"Pan: {abs(pan_pct):.0f}% {pan_dir}  Tilt: {abs(tilt_pct):.0f}% {tilt_dir}  [{servo_status}]"
            cv2.putText(frame, info_text, (10, 30), cv2.FONT_HERSHEY_SIMPLEX, 0.8, (0, 255, 255), 2)
            cv2.putText(frame, f"Persons: {n_persons}  Groups: {len([c for c in clusters.values() if len(c)>=2])}  [{lock_status}] Press 's' to retarget", 
                       (10, 60), cv2.FONT_HERSHEY_SIMPLEX, 0.6, (0, 255, 0), 2)
        else:
            # No target group - tell servos to hold position
            if servo_controller is not None:
                servo_controller.update(0, 0, has_detection=False)
            servo_status = "HOLDING" if servo_controller and servo_controller.enabled else "NO SERVO"
            lock_status = "LOST" if user_data.lock_active else "NO TARGET"
            cv2.putText(frame, f"Persons: {n_persons} [{lock_status}] [{servo_status}] Press 's' to retarget", (10, 30), 
                       cv2.FONT_HERSHEY_SIMPLEX, 0.7, (0, 255, 255), 2)
        
        # Check for key press (non-blocking)
        # Use full key code for arrow keys (don't mask with 0xFF)
        key = cv2.waitKey(1)
        key_masked = key & 0xFF
        
        if key_masked == ord('s') or key_masked == ord('S'):
            # Retarget to largest group
            user_data.lock_active = False
            user_data.locked_group_center = None
            if largest_id >= 0 and largest_id in group_centers:
                user_data.lock_active = True
                user_data.locked_group_center = (group_centers[largest_id][0], group_centers[largest_id][1])
                print(f"Retargeted to largest group (size: {group_centers[largest_id][2]})")
        # Focus control: Arrow keys OR w/e keys
        # Arrow key codes vary by platform - common values:
        # Linux GTK: Up=65362, Down=65364
        # Linux Qt: Up=16777235, Down=16777237
        # Also accept 'w' for up, 'e' for down as reliable fallback
        elif key == 65362 or key == 16777235 or key_masked == ord('w') or key_masked == ord('W'):
            if servo_controller is not None:
                servo_controller.adjust_focus(+1)
        elif key == 65364 or key == 16777237 or key_masked == ord('x') or key_masked == ord('X'):
            if servo_controller is not None:
                servo_controller.adjust_focus(-1)
        
        # Draw focus slider bar on frame
        if servo_controller is not None and servo_controller.enabled:
            draw_focus_slider(frame, servo_controller.focus_angle, height)
        
        # Convert and set frame
        frame = cv2.cvtColor(frame, cv2.COLOR_RGB2BGR)
        user_data.set_frame(frame)
    
    return Gst.PadProbeReturn.OK

# -----------------------------------------------------------------------------------------------
# Keyboard listener for terminal input (works without GUI focus)
# -----------------------------------------------------------------------------------------------
class KeyboardListener:
    """Listen for keyboard input from terminal in a separate thread"""
    def __init__(self, servo_ctrl):
        self.servo_ctrl = servo_ctrl
        self.running = False
        self.thread = None
    
    def start(self):
        """Start the keyboard listener thread"""
        self.running = True
        self.thread = threading.Thread(target=self._listen, daemon=True)
        self.thread.start()
        print("\n[KEYBOARD] Listening for keys: W=focus up, X=focus down, S=retarget, Q=quit")
    
    def stop(self):
        """Stop the keyboard listener"""
        self.running = False
    
    def _listen(self):
        """Listen for keyboard input"""
        import termios
        import tty
        
        # Save terminal settings
        old_settings = termios.tcgetattr(sys.stdin)
        
        try:
            # Set terminal to raw mode
            tty.setcbreak(sys.stdin.fileno())
            
            while self.running:
                # Check if input is available (non-blocking)
                if select.select([sys.stdin], [], [], 0.1)[0]:
                    char = sys.stdin.read(1)
                    
                    if char.lower() == 'w':
                        if self.servo_ctrl:
                            self.servo_ctrl.adjust_focus(+1)
                    elif char.lower() == 'x':
                        if self.servo_ctrl:
                            self.servo_ctrl.adjust_focus(-1)
                    elif char.lower() == 's':
                        print("[KEY] Retarget requested (press in video window)")
                    elif char.lower() == 'q':
                        print("[KEY] Quit requested")
                        os._exit(0)
        except Exception as e:
            print(f"[KEYBOARD] Error: {e}")
        finally:
            # Restore terminal settings
            termios.tcsetattr(sys.stdin, termios.TCSADRAIN, old_settings)

# Global keyboard listener
keyboard_listener = None

# -----------------------------------------------------------------------------------------------
# Argument parsing
# -----------------------------------------------------------------------------------------------
def parse_args():
    parser = argparse.ArgumentParser(description="Human Group Tracking")
    parser.add_argument('--clustering-eps', type=float, default=100,
                        help='Max distance (pixels) between persons to form a group (default: 100)')
    parser.add_argument('--no-servo', action='store_true',
                        help='Disable servo control')
    parser.add_argument('--pan-channel', type=int, default=0,
                        help='PCA9685 channel for pan servo (default: 0)')
    parser.add_argument('--tilt-channel', type=int, default=1,
                        help='PCA9685 channel for tilt servo (default: 1)')
    parser.add_argument('--focus-channel', type=int, default=2,
                        help='PCA9685 channel for focus servo (default: 2)')
    parser.add_argument('--servo-speed', type=float, default=0.5,
                        help='Servo movement speed 0-1 (default: 0.5)')
    parser.add_argument('--invert-pan', action='store_true',
                        help='Invert pan servo direction')
    parser.add_argument('--invert-tilt', action='store_true',
                        help='Invert tilt servo direction')
    parser.add_argument('--swap-servos', action='store_true',
                        help='Swap pan and tilt servo channels (if wired backwards)')
    args, remaining = parser.parse_known_args()
    return args, remaining

# -----------------------------------------------------------------------------------------------
# Main
# -----------------------------------------------------------------------------------------------
if __name__ == "__main__":
    # Set env
    project_root = Path(__file__).resolve().parent.parent
    os.environ["HAILO_ENV_FILE"] = str(project_root / ".env")
    
    # Parse our args
    args, remaining = parse_args()
    
    # Handle servo swap
    pan_ch = args.pan_channel
    tilt_ch = args.tilt_channel
    if args.swap_servos:
        pan_ch, tilt_ch = tilt_ch, pan_ch
    
    print("=" * 60)
    print("Human Group Tracking")
    print(f"Clustering distance: {args.clustering_eps} pixels")
    if not args.no_servo:
        print(f"PCA9685 Servo control: Pan=CH{pan_ch}, Tilt=CH{tilt_ch}, Focus=CH{args.focus_channel}")
        if args.swap_servos:
            print("  (Pan/Tilt channels SWAPPED from default)")
        print(f"Servo speed: {args.servo_speed}")
        if args.invert_pan:
            print("Pan direction: INVERTED")
        if args.invert_tilt:
            print("Tilt direction: INVERTED")
        print("Focus: Use UP/DOWN arrow keys to adjust")
    else:
        print("Servo control: DISABLED")
    print("=" * 60)
    
    # Initialize servo controller (global)
    if not args.no_servo:
        servo_controller = ServoController(
            pan_channel=pan_ch,
            tilt_channel=tilt_ch,
            focus_channel=args.focus_channel,
            speed=args.servo_speed,
            invert_pan=args.invert_pan,
            invert_tilt=args.invert_tilt
        )
    else:
        servo_controller = None
    
    # Start keyboard listener for focus control
    keyboard_listener = None
    if servo_controller is not None:
        keyboard_listener = KeyboardListener(servo_controller)
        keyboard_listener.start()
    
    # Rebuild sys.argv for GStreamerDetectionApp
    sys.argv = [sys.argv[0]] + remaining
    
    # Add --use-frame for visual overlays (required for drawing)
    if '--use-frame' not in sys.argv:
        sys.argv.append('--use-frame')
    
    # Add custom HEF if not specified
    if '--hef-path' not in sys.argv:
        sys.argv.extend(['--hef-path', DEFAULT_CUSTOM_HEF])
        print(f"Using model: {DEFAULT_CUSTOM_HEF}")
    
    # Add labels JSON if not specified
    if '--labels-json' not in sys.argv and Path(DEFAULT_LABELS_JSON).exists():
        sys.argv.extend(['--labels-json', DEFAULT_LABELS_JSON])
    
    # Create callback class
    user_data = user_app_callback_class(eps=args.clustering_eps)
    
    # Run
    print("\nStarting... (Ctrl+C to stop)\n")
    try:
        app = GStreamerDetectionApp(app_callback, user_data)
        app.run()
    finally:
        # Stop keyboard listener
        if keyboard_listener is not None:
            keyboard_listener.stop()
        # Cleanup servo controller
        if servo_controller is not None:
            print("\nCentering servos and cleaning up...")
            servo_controller.cleanup()

\end{lstlisting}
\end{document}
